% =========================================================================
\section{Ejercicio 1: función de transferencia}\label{sec:transferencia}
% =========================================================================
\begin{prelab}[Función de transferencia]
Obtener la función de transferencia $T(s)=\Vout/\Vin$ del circuito de la \cref{fig:filtro} y demostrar que
adopta la forma
\[
  T(s)=\frac{\Vout}{\Vin}=\frac{T_0}{\omega_0^{2}+\dfrac{\omega_0\,s}{Q}+s^{2}} ,
\]
asumiendo que el amplificador operacional es ideal. Obtener las expresiones de $T_0$, $\omega_0$ y $Q$ en función
de $R_1$, $C_1$, $R_2$, $C_2$, $R_A$ y $R_B$.
\end{prelab}

% Figura: filtro activo paso bajo de segundo orden (figura 1 del guion), con nudos y corrientes.
\begin{figure}[ht]
\centering
\begin{circuitikz}[circuito]
  \draw (0,0) node[op amp, ao, noinv input up] (A) {};
  % salida
  \draw (A.out) -- ++(0.9,0) coordinate (O);
  \draw (O) to[short, *-o] ++(0.8,0) node[right, tens] {$\Vout$};
  % entrada no inversora: nudo P, C2 a tierra
  \draw (A.+) -- ++(-1.5,0) coordinate (P);
  \draw (P) to[C, l_=$C_2$, i>_=$I_{C2}$] ++(0,-2.4) node[ground] {};
  % R2 entre X y P
  \draw (P) to[R, l_=$R_2$, i<_=$I_2$] ++(-2.2,0) coordinate (X);
  % R1 entre la entrada y X
  \draw (X) to[R, l_=$R_1$, i<^=$I_1$] ++(-2.3,0) coordinate (E);
  \draw (E) to[short, -o] ++(-0.2,0) coordinate (Ein) node[above, tens] {$\Vin$};
  \draw (Ein) -- ++(-0.7,0) to[sV, l_=$v_{\mathrm{in}}$] ++(0,-2.6) node[ground] {};
  % C1 entre X y la salida
  \draw (X) -- ++(0,1.9) coordinate (T) to[C, l=$C_1$, i>^=$I_{C1}$] (T -| O) -- (O);
  % red RA-RB en la entrada inversora
  \draw (A.-) -- ++(-0.4,0) -- ++(0,-1.3) coordinate (N);
  \draw (N) to[R, l_=$R_A$] (N -| O) -- (O);
  \draw (N) to[R, l=$R_B$] ++(0,-1.6) node[ground] {};
  % nudos y tensiones
  \fill (X) circle (1.6pt); \fill (P) circle (1.6pt); \fill (N) circle (1.6pt);
  \node[tens, anchor=north west, font=\scriptsize] at ($(X)+(0.06,-0.06)$) {$V_x$};
  \node[tens, anchor=south west, font=\scriptsize] at ($(P)+(0.06,0.06)$) {$V_+$};
  \node[tens, anchor=north east, font=\scriptsize] at ($(N)+(-0.06,-0.06)$) {$V_-$};
\end{circuitikz}
\caption{Filtro activo paso bajo de segundo orden (figura~1 del guion), con la alimentación omitida.
Se indican los nudos del análisis ($V_x$, $V_+$, $V_-$) y los sentidos de corriente elegidos.}
\label{fig:filtro}
\end{figure}


\begin{resolucion}
Trabajamos en el dominio de Laplace con condiciones iniciales nulas, de modo que los condensadores tienen
impedancias $1/(sC_1)$ y $1/(sC_2)$, y aplicamos al AO ideal las reglas (R1) y (R2); la realimentación negativa
la proporciona la red $R_A$--$R_B$. Planteamos la ley de Kirchhoff de corrientes en los nudos $V_-$, $V_+$ y
$V_x$, con los sentidos de la \cref{fig:filtro}.

\textbf{1. Nudo inversor: divisor $R_A$--$R_B$.} Como $I_-=0$ (R1), por $R_A$ y $R_B$ circula la misma
corriente, y forman un divisor de tensión de la salida:
\begin{equation}\label{eq:vm}
  V_-(s)=\frac{R_B}{R_A+R_B}\,\Vout(s) .
\end{equation}
Por (R2), $V_+=V_-$. Definiendo la ganancia del amplificador no inversor,
\begin{equation}\label{eq:K}
  K\equiv1+\frac{R_A}{R_B}=\frac{R_A+R_B}{R_B},
\end{equation}
la tensión en la entrada no inversora queda ligada a la salida:
\begin{equation}\label{eq:vp}
  V_+(s)=V_-(s)=\frac{\Vout(s)}{K} .
\end{equation}

\textbf{2. Nudo $V_+$.} Como $I_+=0$ (R1), toda la corriente que llega por $R_2$ se va por $C_2$ a tierra:
\begin{equation}\label{eq:nudoP}
  I_2=I_{C2}\quad\Longrightarrow\quad\frac{V_x-V_+}{R_2}=sC_2\,V_+ .
\end{equation}
Despejamos $V_x$ y sustituimos \eqref{eq:vp}:
\begin{equation}\label{eq:vx}
  V_x=(1+sR_2C_2)\,V_+=\frac{1+sR_2C_2}{K}\,\Vout .
\end{equation}

\textbf{3. Nudo $V_x$.} La corriente que entra por $R_1$ sale por $R_2$ y por $C_1$ (hacia la salida):
\begin{equation}\label{eq:nudoX}
  I_1=I_2+I_{C1}\quad\Longrightarrow\quad\frac{\Vin-V_x}{R_1}=\frac{V_x-V_+}{R_2}+sC_1\,(V_x-\Vout) .
\end{equation}
Por \eqref{eq:nudoP} y \eqref{eq:vp}, el primer sumando del segundo miembro vale $sC_2V_+=sC_2\Vout/K$.
Multiplicando por $R_1$ y despejando $\Vin$:
\begin{equation}\label{eq:vin1}
  \Vin=(1+sR_1C_1)\,V_x+\frac{sR_1C_2}{K}\,\Vout-sR_1C_1\,\Vout .
\end{equation}

\textbf{4. Sustitución.} Llevamos \eqref{eq:vx} a \eqref{eq:vin1} y sacamos factor común $\Vout/K$:
\begin{equation}\label{eq:vin2}
  \Vin=\frac{\Vout}{K}\Bigl[(1+sR_1C_1)(1+sR_2C_2)+sR_1C_2-sK R_1C_1\Bigr] .
\end{equation}
Desarrollando el producto, $(1+sR_1C_1)(1+sR_2C_2)=1+s(R_1C_1+R_2C_2)+s^{2}R_1R_2C_1C_2$, y agrupando por
potencias de $s$:
\begin{equation}\label{eq:vin3}
  \Vin=\frac{\Vout}{K}\Bigl[s^{2}R_1R_2C_1C_2+s\bigl(R_1C_2+R_2C_2+(1-K)R_1C_1\bigr)+1\Bigr] .
\end{equation}

\textbf{5. Función de transferencia.} Dividiendo $\Vout$ entre $\Vin$ en \eqref{eq:vin3}:
\begin{equation}\label{eq:T1}
  T(s)=\frac{\Vout}{\Vin}=\frac{K}{s^{2}R_1R_2C_1C_2+s\bigl(R_1C_2+R_2C_2+(1-K)R_1C_1\bigr)+1} .
\end{equation}
Para llevarla a la forma pedida, dividimos numerador y denominador por $R_1R_2C_1C_2$, de modo que el
coeficiente de $s^{2}$ sea $1$:
\begin{equation}\label{eq:T2}
  \resultado{T(s)=\frac{\dfrac{K}{R_1R_2C_1C_2}}{s^{2}+s\Bigl(\dfrac{1}{R_1C_1}+\dfrac{1}{R_2C_1}+\dfrac{1-K}{R_2C_2}\Bigr)+\dfrac{1}{R_1R_2C_1C_2}}}
\end{equation}
que tiene exactamente la forma del enunciado: una constante en el numerador y un polinomio de segundo grado
en $s$ en el denominador. Es, por tanto, un filtro paso bajo de segundo orden.

\textbf{6. Identificación de los parámetros.} Comparando \eqref{eq:T2} término a término con la forma
del enunciado, $T_0/(s^{2}+\omega_0 s/Q+\omega_0^{2})$:
\begin{align}
  \omega_0^{2}&=\frac{1}{R_1R_2C_1C_2}
  &&\Longrightarrow\quad \resultado{\omega_0=\frac{1}{\sqrt{R_1R_2C_1C_2}}}\label{eq:w0}\\
  T_0&=\frac{K}{R_1R_2C_1C_2}=K\,\omega_0^{2}
  &&\Longrightarrow\quad \resultado{T_0=\frac{1+R_A/R_B}{R_1R_2C_1C_2}}\label{eq:T0}\\
  \frac{\omega_0}{Q}&=\frac{1}{R_1C_1}+\frac{1}{R_2C_1}+\frac{1-K}{R_2C_2}. &&\label{eq:w0Q}
\end{align}
Despejando $Q$ de \eqref{eq:w0Q} con \eqref{eq:w0}, y multiplicando numerador y denominador por
$R_1R_2C_1C_2$:
\begin{equation}\label{eq:Q}
  Q=\frac{\omega_0}{\dfrac{1}{R_1C_1}+\dfrac{1}{R_2C_1}+\dfrac{1-K}{R_2C_2}}
   =\frac{\sqrt{R_1R_2C_1C_2}}{R_1C_2+R_2C_2+(1-K)\,R_1C_1} ,
\end{equation}
y, como $1-K=-R_A/R_B$ por \eqref{eq:K},
\begin{equation}\label{eq:Qfinal}
  \resultado{Q=\frac{\sqrt{R_1R_2C_1C_2}}{(R_1+R_2)\,C_2-\dfrac{R_A}{R_B}\,R_1C_1}}
\end{equation}

\textbf{7. Comprobaciones y comentarios.}
\begin{itemize}
  \item \emph{Límites físicos.} Para $s\to0$ los condensadores son circuitos abiertos: no circula corriente por
        $R_1$ ni $R_2$, $V_+=\Vin$ y el circuito es un no inversor, $\Vout=K\Vin$. En efecto, \eqref{eq:T1} da
        $T(0)=K$. Para $s\to\infty$, $C_2$ cortocircuita $V_+$ a tierra y $\Vout\to0$; en efecto,
        $T(s)\simeq T_0/s^{2}\to0$.
  \item \emph{Papel de cada componente.} $\omega_0$ \eqref{eq:w0} solo depende de las resistencias y
        condensadores de la red RC; la ganancia $K$ solo entra en $T_0$ \eqref{eq:T0} y en $Q$ \eqref{eq:Qfinal}.
        Aumentar $R_A/R_B$ reduce el denominador de \eqref{eq:Qfinal} y, por tanto, \emph{aumenta} $Q$: $C_1$
        devuelve parte de la salida al nudo $V_x$ (realimentación positiva), que realza la respuesta cerca de
        $\omega_0$.
  \item \emph{Estabilidad.} Los polos tienen parte real negativa si $Q>0$, es decir, si el denominador de
        \eqref{eq:Qfinal} es positivo: $R_A/R_B<(R_1+R_2)C_2/(R_1C_1)$. Si se supera ese valor, el circuito
        deja de ser un filtro y oscila.
\end{itemize}
\end{resolucion}
